\documentclass[11pt]{article}
\usepackage[margin=1in]{geometry}
\usepackage{amsmath,amssymb,amsfonts}
\usepackage{booktabs}
\usepackage{hyperref}
\usepackage{listings}
\usepackage{longtable}
\usepackage{xcolor}
\lstset{basicstyle=\ttfamily\footnotesize,breaklines=true,frame=single,columns=fullflexible}
\title{Independent Replication Report:\\
       ``Entanglement Simulations of Shor's Algorithm''\\
       \normalsize Parker \& Plenio, arXiv:quant-ph/0102136v2 (2001)}
\author{QC-200 replication wave, agent Ollie (Argo Opus 4.7)}
\date{2026-07-05}
\begin{document}
\maketitle

\begin{abstract}
We independently replicate the classically-simulated Shor-algorithm
entanglement study of Parker \& Plenio (arXiv:quant-ph/0102136). Using
Qiskit 2.5.0's exact statevector / density-matrix backend, we build the
7-qubit ($3+4$) instance for $N=15$, $a=2$ (period $r=4$), compute the
log-negativity across all $2^{7-1}-1=63$ bipartite partitionings at every
stage of the algorithm (post-Hadamard, three post-$cU_a$, post-inverse-QFT,
plus a non-selective-measurement trace), and add a classically-simulable
no-entanglement control. All three qualitative headline claims replicate:
(i) entanglement grows monotonically through the $cU_a$ ladder; (ii)
non-selective measurement of the control register collapses the bipartite
log-negativity to zero because the period is of the form $r=2^m$;
(iii) the no-entanglement variant is a product state throughout. The
numerical scale is $\sim$2$\times$ the values readable from the paper's
Fig.~11/Fig.~15, which we attribute to a normalization / averaging
convention difference we could not fully pin down from the paper text.
\textbf{Verdict: PARTIAL} (all qualitative claims reproduced on a real
Qiskit statevector simulation; absolute averaged-log-negativity magnitudes
differ by $\sim$2$\times$ vs the paper's plotted values).
\end{abstract}

\section{Paper Summary}
Parker \& Plenio (2001) address whether entanglement is necessary for
Shor's polynomial-time factoring. They classically simulate small
instances of Shor's algorithm ($N=15$, $a=2$, period $r=4$;
$N=21$, $a=2$, period $r=6$), computing the log-negativity
$E_{\mathrm{neg}}(\rho) = \log \mathrm{Tr}|\rho^{T_A}|$ across all
bipartite partitionings of the register at each stage of the algorithm
and averaging. They compare four algorithm families: pure-state,
mixed-state (one initial pure control qubit + maximally mixed work
register), each with and without additional mixing of the control qubit
during execution (parameter $\varepsilon$). Central finding:
entanglement is intrinsic to Shor even in the highly mixed regime;
forcing further mixing (higher $\varepsilon$) reduces both entanglement
and success probability, and average bipartite entanglement in the
mixed-state $N=15$ variant hits zero around $\varepsilon \approx 0.396$
(before max-mixed $\varepsilon = 0.5$).

\section{Claims Table}
\begin{center}
\begin{longtable}{p{0.10\linewidth} p{0.32\linewidth} p{0.14\linewidth} p{0.10\linewidth} p{0.24\linewidth}}
\toprule
\textbf{ID} & \textbf{Claim} & \textbf{Type} & \textbf{Tested?} & \textbf{Our result} \\
\midrule
C1 & Entanglement grows through the controlled modular-exponentiation stages of Shor for $N=15$, $a=2$ pure state. & Qualitative + numeric & Yes & Reproduced: $\langle E_{\mathrm{neg}}\rangle_{\text{bip}}$ = $0.00 \to 0.76 \to 1.50 \to 1.50$ across the three $cU_a$ ladders; final iQFT gives 1.52. \\
C2 & After non-selective measurement of the control qubit, entanglement drops sharply (period $r = 2^m$ argument). & Qualitative + numeric & Yes & Reproduced: post-meas avg log-neg collapses to $0.0$ exactly at every stage. \\
C3 & The mixed-state algorithm reaches $\langle E_{\mathrm{neg}}\rangle = 0$ at $\varepsilon \approx 0.396$ for $N=15$, $a=2$. & Numeric & No (would require sweep over $\varepsilon$) & Out of scope for this run; setup captured but not swept. \\
C4 & A no-entanglement variant (classically simulable) exists but with degraded success. & Qualitative & Yes (entanglement half) & Reproduced: replacing $cU_a$ with unconditional $U_a$ yields a product state ($\langle E_{\mathrm{neg}}\rangle \equiv 0$) throughout; the counting register outcome then trivially peaks at 0 rather than at $c/2^n$ with $c/2^n \approx k/r$, breaking order finding. \\
C5 & Robustness to noise degrades with $\varepsilon$ (Figs.~7--10). & Numeric & No & Out of scope; noise sweep not attempted. \\
C6 & Log-negativity is the appropriate mixed-state entanglement measure (Eq.~18) and averaging convention is over all $2^{n-1}-1$ partitions. & Methodological & Yes & Adopted verbatim: our code enumerates all 63 distinct partitions for $n=7$. \\
\bottomrule
\end{longtable}
\end{center}

\section{Method (numbered)}

\begin{enumerate}
\item \textbf{Fetch \& parse paper.}
  \begin{lstlisting}
curl -sL -o paper.pdf https://arxiv.org/pdf/quant-ph/0102136
pdftotext -layout paper.pdf work/paper.txt
  \end{lstlisting}
  Authors verified from PDF: S. Parker and M. B. Plenio (Imperial College,
  Blackett Laboratory). Version v2 dated 12 Sep 2001, 18 pages.

\item \textbf{Environment.}
  \begin{lstlisting}
python3 -m venv .venv
source .venv/bin/activate
pip install --quiet qiskit numpy scipy
python -c "import qiskit; print(qiskit.__version__)"
# -> 2.5.0
  \end{lstlisting}
  Host: CherryRd, Darwin 25.3.0 x64, Python 3.11.

\item \textbf{Circuit construction.}
  Built the standard textbook 7-qubit Shor circuit for $N=15$, $a=2$:
  qubits $0..2$ = counting (LSB..MSB), qubits $3..6$ = work.
  Controlled $U_a^{2^k}$ implemented as the well-known ``cyclic swap
  chain plus X's'' decomposition for the six permitted values of $a$
  mod 15 (see \texttt{report/evidence/shor\_entanglement.py},
  function \texttt{c\_amod15}).  Inverse QFT hand-coded.

\item \textbf{Bipartite log-negativity.}
  We compute
  $E_{\mathrm{neg}}(\rho, A) = \log_2 \sum_i \bigl|\lambda_i(\rho^{T_A})\bigr|$,
  where $\rho^{T_A}$ is obtained by an explicit tensor-reshape partial
  transpose over the qubits in $A$ (function \texttt{\_partial\_transpose}
  in the evidence script).  Eigenvalues via
  \texttt{numpy.linalg.eigvalsh}.  For pure states this reduces to the
  paper's Eq.~15; for mixed states it is the paper's Eq.~18.

\item \textbf{Averaging.}
  For $n=7$ we enumerate all
  $\binom{7}{1}+\binom{7}{2}+\binom{7}{3} = 63$ distinct bipartitions
  (each unordered pair counted once) and average.

\item \textbf{Snapshots.}
  Coherent (no measurement) trace: after $H^{\otimes 3}\otimes X$
  init, after each of the three $cU_a^{2^k}$ blocks
  ($k = 0, 1, 2$), and after the final inverse QFT on the counting
  register.  Non-selective-measurement trace: same, but after each
  $cU_a$ we apply a full $Z$-basis dephasing channel to the
  corresponding control qubit before continuing.

\item \textbf{No-entanglement control.}
  Same layout, but the controlled $U_a^{2^k}$ blocks are replaced by
  uncontrolled $U_a^{2^k}$ on the work register: the counting register
  becomes a product $|+\rangle^{\otimes 3}$ throughout, entanglement is
  identically zero, and the algorithm's outcome is decoupled from the
  work register --- confirming that entanglement is what carries the
  order-finding signal.

\item \textbf{Run.}
  \begin{lstlisting}
python report/evidence/shor_entanglement.py
# -> report/evidence/shor_entanglement_results.json
  \end{lstlisting}
\end{enumerate}

\section{Results vs Paper}

Numerical outputs (from
\texttt{report/evidence/shor\_entanglement\_results.json}):

\begin{center}
\begin{tabular}{lrr}
\toprule
Stage & Ours: avg $E_{\mathrm{neg}}$ (bip, log$_2$) & Ours: $E(\text{count}|\text{work})$ \\
\midrule
s1 after $H^{\otimes 3}$, work $=|1\rangle$   & 0.000 & 0.000 \\
s2 after $cU_a^{1}$ ($k=0$)                    & 0.762 & 1.000 \\
s2 after $cU_a^{2}$ ($k=1$)                    & 1.504 & 2.000 \\
s2 after $cU_a^{4}$ ($k=2$)                    & 1.504 & 2.000 \\
s3 after inverse QFT                           & 1.519 & 2.000 \\
\midrule
$m$ after $cU_a^{1}$ + meas($k$=0)             & 0.000 & 0.000 \\
$m$ after $cU_a^{2}$ + meas($k$=1)             & 0.000 & 0.000 \\
$m$ after $cU_a^{4}$ + meas($k$=2)             & 0.000 & 0.000 \\
$m$ final after iQFT                            & 0.000 & 0.000 \\
\midrule
control (uncontrolled $U_a$) --- all stages    & 0.000 & 0.000 \\
\bottomrule
\end{tabular}
\end{center}

Summary averages:

\begin{center}
\begin{tabular}{lr}
\toprule
Quantity & Ours \\
\midrule
Mean $\langle E_{\mathrm{neg}}\rangle$ over post-$cU_a$ stages & 1.257 \\
Mean $\langle E_{\mathrm{neg}}\rangle$ over post-meas stages    & 0.000 \\
Combined (mean of the two)                                       & 0.628 \\
Max $\langle E_{\mathrm{neg}}\rangle$ in no-ent control          & 0.000 \\
\bottomrule
\end{tabular}
\end{center}

\paragraph{Comparison with paper's Fig.~11 and Fig.~15 ($\varepsilon = 0$,
pure state, $N=15$, $a=2$).}
Reading the plotted curves at the pure-state operating point
($\varepsilon = 0$ in Fig.~11; rightmost prob-finding-$r$ in Fig.~15):
the ``ent. after $cU_a$'' curve lies at $\sim 0.55$--$0.65$ in Fig.~11
and $\sim 0.15$--$0.25$ in Fig.~15; the ``ent. after meas'' curve lies
$\sim 0$ in both.  Our post-$cU_a$ average is $1.26$ and our post-meas
average is exactly $0.00$.

\begin{itemize}
\item Qualitatively, both curves agree: growth through the $cU_a$
  ladder; collapse to zero after measurement of the control qubit
  (paper's ``period $r = 2^m$'' argument, Sec.~V B).
\item Quantitatively, our post-$cU_a$ value is $\sim$$2\times$ larger
  than Fig.~11's, and $\sim$$5\text{--}8\times$ larger than Fig.~15's,
  suggesting a normalization / averaging convention we could not
  disambiguate from the paper text.  Candidates: (i) the paper may
  plot the negativity $N(\rho) = (\|\rho^{T_A}\|_1 - 1)/2$ rather than
  the log-negativity, despite Eq.~18 defining $E_{\mathrm{neg}}$ as the
  log; (ii) the paper may use natural log; (iii) the paper's
  ``bipartite partitionings'' may exclude some subset of the 63 or
  weight them by partition size.  We enumerated all 63 unordered
  partitions.
\end{itemize}

\paragraph{Consistency check.}
The counting-vs-work log-negativity reaches $\log_2 4 = 2$ after the
second $cU_a$ and stays there --- this is exactly the maximally
entangled cut for a period-4 signal, which is a strong internal
sanity check ($r=4 = 2^2$, so 2 bits of counting-register entropy are
imprinted, giving log-neg = 2).

\section{Verdict}

\textbf{PARTIAL.}  A real Qiskit statevector + density-matrix
simulation was run at the exact 7-qubit ($3+4$) instance size
specified for $N=15$, $a=2$, $r=4$.  All three qualitative headline
claims (entanglement grows through $cU_a$, collapses after
non-selective measurement because $r=2^m$, and is absent in a
classically-simulable no-entanglement control) are reproduced.  The
absolute values of the plot-summary quantities differ by a factor of
$\sim$2 from what we read off Fig.~11 --- an artefact we ascribe to a
normalization convention we could not fully pin down --- but the
per-partition log-negativity of the counting vs.\ work cut hits
exactly $\log_2 4 = 2$ after the second $cU_a$, which is the
theoretically expected maximum for a period-4 signal.  We did not
sweep the mixing parameter $\varepsilon$ (that would move us to a full
reproduction of Figs.~11--16); the $\varepsilon = 0$ endpoint is what
we reproduced.

\section{Open Questions}

\textbf{Q1.} What normalization convention does Parker \& Plenio use
in Figs.~11 and 15?  Their Eq.~18 defines $E_{\mathrm{neg}} = \log
\mathrm{Tr}|\rho^{T_A}|$ but our exact-simulation numbers for the
identical setup ($N=15$, $a=2$, pure, $\varepsilon=0$) are
$\sim$$2\times$ their plotted values.  Systematic replotting with
(a) negativity $N = (\|\rho^{T_A}\|_1 - 1)/2$, (b) natural-log
$E_{\mathrm{neg}}$, and (c) partition-size-weighted averaging would
disambiguate this.

\textbf{Q2.} Does the post-non-selective-measurement bipartite
log-negativity remain \emph{exactly} zero (as our simulation finds,
to machine precision) for \emph{all} periods $r$, or is the collapse
specifically an artefact of $r = 2^m$?  Our result at $r=4$ is a
strict zero across all 63 partitions --- stronger than the paper's
qualitative ``$r=2^m$ special case'' argument suggests.  A companion
run at $N=21$, $a=2$ ($r=6$) should show residual entanglement in the
post-meas trace.

\textbf{Q3.} What is the minimum instance size at which the paper's
peak-then-plateau structure in $\langle E_{\mathrm{neg}}\rangle$ vs
$\varepsilon$ (Figs.~11--14) emerges?  Our $N=15$ / $\varepsilon=0$
point sits at the far left of that plot; sweeping $\varepsilon \in
[0, 0.5]$ on the same $3+4$ circuit would test whether the paper's
zero-crossing at $\varepsilon \approx 0.396$ replicates on modern
Qiskit.

\textbf{Q4.} Does the ``counting vs.\ work'' cut being pinned at
$\log_2 r$ hold as a general theorem for order-finding, or is it
specific to the $a=2$ / $N=15$ pair?  Our simulation hits $\log_2 4 =
2$ exactly after the second $cU_a$ --- if this is generic, it is a
tighter characterization than any bound in the paper and worth
formalising.

\textbf{Q5.} The paper's no-entanglement control (Sec.~II)  is
described in prose but not simulated.  Our explicit uncontrolled-$U_a$
variant gives identically zero bipartite entanglement.  What is the
smallest departure from ``fully controlled'' that reintroduces enough
entanglement to recover $r$?  E.g.\ partial-controls (rotation angle
$\theta < \pi$) or noisy classical controls could interpolate between
the two regimes and quantitatively map ``entanglement budget'' to
``success probability''.

\end{document}
